% !PRINTED 315
\weberpage{315}
%<<CONT
regions, then, the loss of changes of sign gives the correct number of
roots. In the region 2, on the other hand, the decision between two
roots and none is not made.

Let us now take the \gesperrt{Jacobi} criterion (\S.\,102). According
to it, in order to estimate the number of roots lying between $\alpha$
and $\beta$, we must first transform the equation (1) into
\[
(\alpha+\beta y)^2+a(\alpha+\beta y)(1+y)+b(1+y)^2=0
\]
or
\[
\begin{aligned}
&(\beta^2+a\beta+b)y^2+(2\alpha\beta+a(\alpha+\beta)+2b)y\\
&\quad{}+(\alpha^2+a\alpha+b)=0,
\end{aligned}
\]
and now count the number of changes of sign in the three functions
\[
\begin{aligned}
&\beta^2+a\beta+b\\
&2\alpha\beta+a(\alpha+\beta)+2b\\
&\alpha^2+a\alpha+b
\end{aligned}
\]
The first and the third of these expressions, set equal to zero,
represent the two parabola tangents considered above. The middle one,
$2\alpha\beta+a(\alpha+\beta)+2b$, vanishes for $a=-2\alpha$,
$b=\alpha^2$ and for $a=-2\beta$, $b=\beta^2$ and thus represents the
line joining the two points $\alpha$, $\beta$, and indeed in such a way
that it becomes negative on the side that contains the point of
intersection of the two tangents [at this point of intersection itself
it is equal to $-(\beta-\alpha)^2$].

% SOURCE FIGURE CROP: Bd. I registered scan 341 (printed p. 315), 600 dpi,
% crop x=100,y=3000,w=700,h=620; source-faithful raster, Fig. 20.
\weberfigureimage{weber-b1-fig20.png}{0.34}
We obtain here five regions with no, one or two changes of sign, as
Fig. 20 indicates.

One sees from this that this figure confines the indecision to a much
smaller region than the previous one. From this point of view, then, it
does not appear justified that \gesperrt{Jacobi} so decidedly gives the
\gesperrt{Budan-Fourier} criterion preference over his
own\footnote{Cf. F.\ \gesperrt{Klein}, Geometrisches zur
Abz\"ahlung der reellen Wurzeln einer algebraischen Gleichung in
\gesperrt{Dyck}'s Catalog mathematischer Modelle (M\"unchen 1892).}.
